% GENERATED FILE. Edit canonical lessons, then run npm run generate:books.
% canonical-source-hash: fnv1a64:0bed9a577c775e77
% canonical-lessons: PA-C135-gardan, PA-C135-chati, PA-C135-pitth, PA-C135-lakk, PA-C135-kuhni

\chapter{Neck, Chest, Back, Waist, Elbow}
\label{ch:pa-neck-chest-back-waist-elbow}
\hlchaptermodality{135}

\begin{chapteropening}
\textbf{By the end of this chapter:} \emph{I can say a neck, a chest, a back, a waist and an elbow.}

Complete the last lesson of chapter 135: I can say a neck, a chest, a back, a waist and an elbow.
\end{chapteropening}

\section[\texorpdfstring{\pa{ਗਰਦਨ} (gardan)}{ਗਰਦਨ (gardan)}]{\pa{ਗਰਦਨ} (gardan) — a neck}

\label{lesson:PA-C135-gardan}



\begin{quote}
Before the new one: say the Punjabi for a beard, then the Punjabi for a moustache.
\end{quote}

\subsection*{You'll want to know: \pa{ਗਰਦਨ}}
\textbf{\pa{ਗਰਦਨ}} — \emph{gardan} — \textquotedblleft{}a neck\textquotedblright{}.

\begin{cousinweb}
\textbf{\pa{ਗਰਦਨ}} (\emph{gardan}) is a loan, usually traced to Persian \emph{gardan}, \textquotedblleft{}neck\textquotedblright{}.
\end{cousinweb}

\begin{grammarlens}[title={What you've built}]
One more word for this chapter.
\end{grammarlens}

\subsection*{Your turn}
\begin{itemize}
\raggedright
  \item \emph{Say it:} \emph{gardan}
  \item \emph{Say it:} \emph{gardan}, once more
  \item \emph{Say it:} say \emph{gardan}
  \item \emph{Recall:} say the Punjabi for a beard, then the Punjabi for a moustache, then say \emph{gardan} again
\end{itemize}

\begin{tcolorbox}[breakable,colback=teal!4,colframe=teal!35!black,title={Before you move on}]
What does \pa{ਗਰਦਨ} mean? (\textquotedblleft{}A neck\textquotedblright{}.) Say it once more.
\end{tcolorbox}
\section[\texorpdfstring{\pa{ਛਾਤੀ} (chātī)}{ਛਾਤੀ (chātī)}]{\pa{ਛਾਤੀ} (chātī) — a chest}

\label{lesson:PA-C135-chati}



\begin{quote}
Before the new one: say the Punjabi for a moustache, then the Punjabi for a neck.
\end{quote}

\subsection*{You'll want to know: \pa{ਛਾਤੀ}}
\textbf{\pa{ਛਾਤੀ}} — \emph{chātī} — \textquotedblleft{}a chest\textquotedblright{}.

\begin{cousinweb}
\textbf{\pa{ਛਾਤੀ}} (\emph{chātī}) is a tadbhava, a word inherited and reshaped by sound change. It comes from \emph{*chātti}, \textquotedblleft{}chest, breast\textquotedblright{}, a form CDIAL reconstructs from its descendants (CDIAL 5014).
\end{cousinweb}

\begin{grammarlens}[title={What you've built}]
One more word for this chapter.
\end{grammarlens}

\subsection*{Your turn}
\begin{itemize}
\raggedright
  \item \emph{Say it:} \emph{chātī}
  \item \emph{Say it:} \emph{chātī}, once more
  \item \emph{Say it:} say \emph{chātī}
  \item \emph{Recall:} say the Punjabi for a moustache, then the Punjabi for a neck, then say \emph{chātī} again
\end{itemize}

\begin{tcolorbox}[breakable,colback=teal!4,colframe=teal!35!black,title={Before you move on}]
What does \pa{ਛਾਤੀ} mean? (\textquotedblleft{}A chest\textquotedblright{}.) Say it once more.
\end{tcolorbox}
\section[\texorpdfstring{\pa{ਪਿੱਠ} (piṭṭh)}{ਪਿੱਠ (piṭṭh)}]{\pa{ਪਿੱਠ} (piṭṭh) — a back}

\label{lesson:PA-C135-pitth}



\begin{quote}
Before the new one: say the Punjabi for a neck, then the Punjabi for a chest.
\end{quote}

\subsection*{You'll want to know: \pa{ਪਿੱਠ}}
\textbf{\pa{ਪਿੱਠ}} — \emph{piṭṭh} — \textquotedblleft{}a back\textquotedblright{}.

\begin{cousinweb}
\textbf{\pa{ਪਿੱਠ}} (\emph{piṭṭh}) is a tadbhava, a word inherited from Sanskrit and reshaped by sound change. It comes from \emph{pṛṣṭi}, \textquotedblleft{}vertebrae; back\textquotedblright{} (CDIAL 8370).
\end{cousinweb}

\begin{grammarlens}[title={What you've built}]
One more word for this chapter.
\end{grammarlens}

\subsection*{Your turn}
\begin{itemize}
\raggedright
  \item \emph{Say it:} \emph{piṭṭh}
  \item \emph{Say it:} \emph{piṭṭh}, once more
  \item \emph{Say it:} say \emph{piṭṭh}
  \item \emph{Recall:} say the Punjabi for a neck, then the Punjabi for a chest, then say \emph{piṭṭh} again
\end{itemize}

\begin{tcolorbox}[breakable,colback=teal!4,colframe=teal!35!black,title={Before you move on}]
What does \pa{ਪਿੱਠ} mean? (\textquotedblleft{}A back\textquotedblright{}.) Say it once more.
\end{tcolorbox}
\section[\texorpdfstring{\pa{ਲੱਕ} (lakk)}{ਲੱਕ (lakk)}]{\pa{ਲੱਕ} (lakk) — a waist}

\label{lesson:PA-C135-lakk}



\begin{quote}
Before the new one: say the Punjabi for a chest, then the Punjabi for a back.
\end{quote}

\subsection*{You'll want to know: \pa{ਲੱਕ}}
\textbf{\pa{ਲੱਕ}} — \emph{lakk} — \textquotedblleft{}a waist\textquotedblright{}.

\begin{cousinweb}
\textbf{\pa{ਲੱਕ}} (\emph{lakk}) is a tadbhava, a word inherited and reshaped by sound change. It comes from \emph{*lakka}, \textquotedblleft{}waist, loins\textquotedblright{}, a form CDIAL reconstructs from its descendants (CDIAL 10879).
\end{cousinweb}

\begin{grammarlens}[title={What you've built}]
One more word for this chapter.
\end{grammarlens}

\subsection*{Your turn}
\begin{itemize}
\raggedright
  \item \emph{Say it:} \emph{lakk}
  \item \emph{Say it:} \emph{lakk}, once more
  \item \emph{Say it:} say \emph{lakk}
  \item \emph{Recall:} say the Punjabi for a chest, then the Punjabi for a back, then say \emph{lakk} again
\end{itemize}

\begin{tcolorbox}[breakable,colback=teal!4,colframe=teal!35!black,title={Before you move on}]
What does \pa{ਲੱਕ} mean? (\textquotedblleft{}A waist\textquotedblright{}.) Say it once more.
\end{tcolorbox}
\section[\texorpdfstring{\pa{ਕੂਹਣੀ} (kūhṇī)}{ਕੂਹਣੀ (kūhṇī)}]{\pa{ਕੂਹਣੀ} (kūhṇī) — an elbow}

\label{lesson:PA-C135-kuhni}



\begin{quote}
Before the new one: say the Punjabi for a back, then the Punjabi for a waist.
\end{quote}

\subsection*{You'll want to know: \pa{ਕੂਹਣੀ}}
\textbf{\pa{ਕੂਹਣੀ}} — \emph{kūhṇī} — \textquotedblleft{}an elbow\textquotedblright{}.

\begin{cousinweb}
\textbf{\pa{ਕੂਹਣੀ}} (\emph{kūhṇī}) is a tadbhava, a word inherited from Sanskrit and reshaped by sound change. It comes from \emph{kaphoṇi}, \textquotedblleft{}elbow\textquotedblright{} (CDIAL 2757). CDIAL takes the Sanskrit word, probably, from Munda rather than Dravidian.
\end{cousinweb}

\begin{grammarlens}[title={What you've built}]
One more word for this chapter.
\end{grammarlens}

\subsection*{Your turn}
\begin{itemize}
\raggedright
  \item \emph{Say it:} \emph{kūhṇī}
  \item \emph{Say it:} \emph{kūhṇī}, once more
  \item \emph{Say it:} say \emph{kūhṇī}
  \item \emph{Recall:} say the Punjabi for a back, then the Punjabi for a waist, then say \emph{kūhṇī} again
\end{itemize}

\begin{tcolorbox}[breakable,colback=teal!4,colframe=teal!35!black,title={Before you move on}]
What does \pa{ਕੂਹਣੀ} mean? (\textquotedblleft{}An elbow\textquotedblright{}.) Say it once more.
\end{tcolorbox}
